We first prove a comparison theorem for local passage laws.  It will be
used for critical FK maps and, later, for spanning-tree-decorated maps.
The two applications construct their conditional laws and initialize their
distance units separately.  The common argument identifies a law once it
has a finite positive optimal upper comparison with the known LQG metric.
Theorem~\ref{common:identification} identifies its infimal internal
passage cost $F_U$ as $C D_h(\cdot,\cdot;U)$, with one deterministic
constant for the law. The proof proceeds through annulus comparison,
equality tangents, and the finite forcing argument.

In this section only, fix
\[
 \sqrt2\le\gamma<2,\qquad p=1/d_\gamma,\qquad
 \xi=\gamma p,\qquad Q=2/\gamma+\gamma/2,
\]
and a number $\lambda>1$ for the geometric sequence of quantum volumes.
The applications use $\lambda=4$ for FK and $\lambda=2$ for spanning
trees.  This number changes volume units only; the Euclidean annulus
radii used below are fixed geometric choices.  All constants may depend
on these fixed parameters.  No coupling or uniformity across parameters
is asserted.

Write $D_U=D_h(\cdot,\cdot;U)$ for the internal reference metric.
A path has \emph{open confinement} in $U$ when its range is a compact
subset of $U$.  The reference metric has the properties in
Theorem~\ref{bg:metric-inputs}.  The metric, locality, Weyl scaling,
uniqueness and confluence inputs cited there hold for every fixed
subcritical $\gamma$, including $\sqrt2$.  For the inward
box-connection proof, the estimate of
\cite[Lemma~3.20, Equations~(3.66)--(3.67)]{DFGPSWeakMetric}
allows $0<\chi<\xi(Q-2)$; this interval remains nonempty at
$\gamma=\sqrt2$.  The same inward chain therefore supplies its face
connections at that value as well.  We retain the actual
joint law of the fields, curve, and internal reference metrics.
In particular, a field is not replaced by its class modulo real additive
constants.  On every fixed compact region used for Gaussian field
modifications, choose one whole-plane Gaussian reference whose pinning
circle lies outside that region and all its buffers.  A change of field
law will always keep the same local conditional passage kernel.

\begin{lemma}[Reference continuity at the cone mark]
\label{common:cone-mark-continuity}
For the circle-average $\gamma$-quantum cone at any fixed parameter
in this section, the reference metric extends to its marked point
$0$ and induces the Euclidean topology there.  In particular,
\[
 \operatorname{diam}_{D_h}(B_\varepsilon(0))\longrightarrow0
 \qquad\text{as }\varepsilon\downarrow0.
\]
\end{lemma}
\begin{proof}
Put $r_n=2^{-n}$, $A_n=\{r_n/2\le|z|\le r_n\}$ and
$A_n^+=\{r_n/4<|z|<2r_n\}$.  For $n\ge2$, the cone field on
these annuli is $G_1-\gamma\log|z|$.  Write
$H(r)=h_r(0)+Q\log r$.  The radial cone definition gives
\[
 H(r_n)=-(Q-\gamma)n\log2+o(n)\quad\text{almost surely}.
\]
Conformal covariance and the reference metric's constant Weyl rule
give
\[
 \operatorname{diam}_{D_h(\cdot,\cdot;A_n^+)}(A_n)
       =e^{\xi H(r_n)}Z_n,
\]
where each $Z_n$ has the same marginal law: the internal diameter
of $\{1/2\le|z|\le1\}$ inside $\{1/4<|z|<2\}$ for the field
$G_1-\gamma\log|z|$.  This variable has a positive moment by the
reference internal-diameter moment bound of
\cite[Theorem~1.8]{DFGPSWeakMetric}.  The logarithmic drift is bounded
on this fixed annulus, so the reference Weyl rule preserves that
moment bound.  No independence among the $Z_n$ is needed.

Since $Q-\gamma>0$, Markov's inequality and Borel--Cantelli bound
$Z_n$ by $e^{\delta n}$ eventually, with
$0<\delta<\xi(Q-\gamma)\log2$.  The displayed internal diameters
are therefore summable.  Adjacent closed annuli meet on a circle;
joining through these circles shows that the ambient diameter of
$B_{r_n}(0)\setminus\{0\}$ is at most twice the tail sum.  All
sequences approaching $0$ consequently define the same reference
completion point, and that tail sum tends to zero.  A positive
crossing lower bound in any fixed compact annulus separates this
completion point from every other point and prevents a sequence
staying outside a Euclidean neighborhood of $0$ from approaching it
in the metric.  This proves the claimed topology and diameter limit.
\end{proof}

Here is the precise interface used by the common proofs.

\begin{definition}[Families of local passage laws]
\label{common:admissible}
A family of local passage laws consists of laws of the reference
continuum input together with countable closed collections of path
certificates and their local restrictions.  It has the following
properties.
\begin{enumerate}[label=(\roman*)]
\item Certificates retain endpoints, compact ranges, every finite
number of ordered subpaths, and nonnegative additive cost bounds.
Finite certificates have actual path realizations with arbitrary open
loss; bounded-cost realizing paths have subsequential certificates.
Individual edge costs tend to zero, and ordered marking retains the
subdivision of any bounded-cost path at successive cost increments
$\delta$ with overshoot tending to zero.  The finite separator constructions
used below are valid for these paths, with actual primal intersections
whenever the separators are primal circuits.
\item The law of a restriction compactly inside an open real-field
patch $U$ is a kernel from the actual real field on $U$, local imaginary
traversals in a possibly larger buffer, and local Poisson area marks.
It may also depend on the finite genuine order of those traversals.
The feasible orders have positive conditional probabilities depending
only on the unclocked imaginary input.  Their reference law is the same
at every quantum index.  An order-free kernel is the special case of
one possible order.  Restrictions in disjoint real-field patches have
product kernels conditional on the complete continuum input.  These
kernels are compatible under restriction and countable exhaustion.
\item The same local kernels are used after fixed locally equivalent
field changes and intrinsic coordinate changes.  Local area marks have
their correct conditional quantum-area intensities; they may be
integrated out using their independent restrictions.  If $S_a$ multiplies
all recorded costs by $a$ and $\mathcal K$ is a local kernel, write
\[
 c_v=\gamma^{-1}\log v,\qquad
 \mathcal K^{(v)}(g,\widehat g,o)
  =(S_{v^{-p}})_\#\mathcal K(g+c_v,\widehat g,o).
\]
At each fixed volume $v$ these indexed laws are limits of the corresponding actual
time-, height-, and cost-transformed discrete experiments.  The affine
and volume identities are the ones stated in
Equations~\eqref{comp:ann-index-definition}--\eqref{comp:ann-two-units}.
\item For each sequence of laws and changes of base units used in the
argument, the model supplies uniform local tightness and modulus bounds
for the coordinates it retains.  These bounds give relative compactness
and a jointly convergent diagonal discrete subsequence.  In it, all marked
observations, reference experiments, fixed quantum offsets, and actual
joint field/internal-metric readouts can be retained under further
extraction.  This requirement includes the varying laws selected in
amplification and equality-tangent extraction, with the bounds uniform
along each selected indexed sequence; it makes no compactness
claim for uncontrolled sequences of laws.  Each fixed retained time unit still tends to
infinity in lattice steps.  Additional continuum coordinates are drawn
from their canonical conditional law independently of the entire raw
actual and reference observation array given the full contour and cuts.
The local kernel identities hold anew in each resulting law, with its
own extracted reference kernels.  Existing uniform comparison bounds
pass through these marked extractions.
\item Whenever a finite positive upper comparison has been obtained,
continuum certificates concatenate with arbitrarily small cost loss and
open confinement, and their infimal internal costs are locally
continuous.  The comparison bound holds uniformly at all fixed quantum
offsets and in further extractions to which it is applied.
\item If there is a nontrivial order input, buffered directed
traversals have the local reconstruction, rational-point visit and
transience properties of Lemma~\ref{comp:local-order-certificate}.
The complete dependency disks of its local determination events can be
used for independent Gaussian trials.  In the one-order case this
condition is vacuous.
\end{enumerate}
All restrictions and full observation spaces are standard Borel.  Thus
prescribed separated restrictions with their product conditional law can
be completed by a regular conditional distribution to a full sample.
\end{definition}

Comparison assertions are tested on a countable exhaustion of ordinary
buffered charts.  When transferring their optimal constants between a
surface law and Gaussian reference laws, the model application must also
verify that its actual coordinate and quantum-unit transformations cover
these charts.  Absolute continuity on one patch alone is not a
whole-surface assertion.  Singular marked endpoints, when present, are
handled by the separate endpoint or completion argument.

These are requirements on the actual coupled laws, not consequences of
contour convergence alone.  The FK kernels of
Theorem~\ref{loc:spatial-kernels} and
Proposition~\ref{loc:kernel-closure} supply the local and extraction
properties.  Their continuum joining property will follow from the
small-circuit lemma below and the FK normalization argument.  The tree
application verifies the same interface using its independent-walk
kernels and its own stronger joining estimate.  In neither application
does this interface assume a Weyl rule for the unidentified passage law.

\subsection{Separating circuits and port splicing}

The certificate operations used in this subsection are simultaneous marked
subpath extraction, the additive cost inequality, realization with open
ports, and robust planar separation.  Each invocation requires these
operations for the actual primal graph paths.  Euclidean intersections of
interpolations alone do not suffice.

\begin{definition}[Upper comparison with ports]\label{comp:port-upper}
An upper comparison $C$ means the following.  If $P$ is a nonconstant
$D$-rectifiable path, $U$ is an open neighborhood of its range, and $O_0,O_1$
are open neighborhoods of its endpoints, there is a passage certificate
confined to $U$, with endpoints in $O_0,O_1$, of cost at most
$C\operatorname{len}_D(P)+\epsilon$, for every $\epsilon>0$.
The assertion is required simultaneously for finite lists of such tests.
A lower comparison $c$ means that every certificate with endpoints $x,y$
and cost $a$ satisfies $a\ge cD(x,y)$.
\end{definition}

Here and subsequently an infimum over certificates with a specified pair of
endpoints refers to the closed certificate collection, not to the limit of
the distance between an arbitrarily prescribed pair of discrete vertices.

\begin{lemma}[Splicing by small circuits]\label{comp:port-splicing}
Suppose that for every $K\Subset U$ and every $\epsilon>0$, every $x\in K$
is surrounded by a primal separating circuit of cost at most $\epsilon$,
whose range is contained in $U\cap B(x,\epsilon)$.  Circuits are required
to be realizable with strict separation from a smaller neighborhood of $x$.
Assume the certificate operations stated above.  Then nonconstant
certificates ending at the same point concatenate with arbitrary cost
loss and open confinement.  If an upper comparison $C<\infty$ also holds,
the infimal cost $F_U$ satisfies
\begin{equation}\label{comp:upper-internal}
 F_U(x,y)\le C D_U(x,y),\qquad x,y\in U,
\end{equation}
and is locally continuous in its two endpoints.  These conclusions do not
bound the cost of escaping from a prescribed discrete vertex.
\end{lemma}

\begin{proof}
Consider certificates from $x$ to $z$ and from $z$ to $y$, with $x,y\ne z$.
Choose a circuit surrounding $z$ that avoids $x$ and $y$ and has cost less
than $\epsilon$.  In sufficiently accurate realizations both paths must
cross this circuit: their endpoints near $z$ lie on its inner side and
their other endpoints on its outer side.  Robust separation makes the
crossings actual common primal vertices.  Cut at such vertices and join
along the circuit.  The resulting cost is at most the sum of the original
costs plus $\epsilon$.  Simultaneous marked extraction retains that
inequality.  Letting the circuit shrink proves concatenation; the cases
with a constant certificate are dealt with by the cost-zero diagonal.

For a path $P$ compactly confined to $U$, apply the upper comparison with
endpoint neighborhoods shrinking to its endpoints.  Extract the resulting
certificates in a fixed compact confinement, first enlarging the range of
$P$ slightly within $U$.  Their endpoints become exact and their costs are
at most $C\operatorname{len}_D(P)$.  Taking the infimum over such paths
proves Equation~\eqref{comp:upper-internal}.  The triangle inequality just
proved and symmetry give
\[
 |F_U(x,y)-F_U(x',y')|
 \le C\{D_U(x,x')+D_U(y,y')\}
\]
whenever these quantities are finite.  On a compact subset of $U$, $D_U$
and $D$ agree for sufficiently close pairs: the compact set has positive
$D$-distance from the complement of $U$, and a sufficiently short
near-minimizing path cannot exit $U$.  Continuity follows.
\end{proof}

\begin{lemma}[Finite separator chains]\label{comp:circle-chain}
Let $K\Subset U$.  For each integer $n$ let $\mathcal B_n$ be a finite
collection of closed Jordan domains of Euclidean diameters at most
$2\epsilon_n$, where $\epsilon_n\downarrow0$, with weights $w_B>0$.
Each domain has a designated open hole compactly in its interior, and
the holes cover $K$.  The domains and the buffered boundary connections
below lie in $U$.

Suppose a continuous path $P$ in $K$ is charged at least $aw_B$ whenever,
starting inside the hole of $B$, it exits $B$.  Suppose also that any two
points of $\partial B$ can be connected along the boundary, or in its
buffered band, with cost at most $Aw_B$.  Then there is a connection between points within
$4\epsilon_n$ of the endpoints of $P$, with cost at most
$(A/a)$ times its total charge.  The connection uses only the buffered
bands of disks selected along $P$.  Its first and last errors are
Euclidean errors.
\end{lemma}

\begin{proof}
At the current point of $P$, choose a disk whose hole contains it and stop
at its first exit.  Repeat, stopping before an incomplete last traversal.
Completed traversal intervals are disjoint.  They are finite in number
when the total charge is finite, since each costs at least
$a\min_{B\in\mathcal B_n}w_B>0$.  Their weights therefore sum to at most
the total charge divided by $a$.

Consecutive selected domains intersect.  Delete domains contained in others,
and replace each domain in the original chain by a maximal selected domain
containing it.  The resulting intersection graph connects a domain
containing the initial point to a domain containing the final completed
exit.  Choose a simple graph path between them.  Two consecutive maximal
Jordan domains have intersecting boundaries: by the Jordan separation
theorem, otherwise one of the two intersecting domains would contain the
other.  Follow the boundary connections between intersection points.  No retained boundary is used
twice, so the cost is at most $A$ times the sum of the selected weights.
Containment deletion moves an endpoint by at most $2\epsilon_n$, and an
incomplete last disk adds at most $2\epsilon_n$.  If no traversal is
completed, both endpoints are in one disk and only this endpoint error
remains.
\end{proof}

The same proof applies on a finite planar map when the domains are the
filled interiors of simple primal circuits and ``diameter'' means the
diameter of their projected images.  The boundary intersections are then
actual primal vertices.  This version is the one used for circuit
quantiles; a cheap surrounding circuit need not pass through every point
of a prescribed Euclidean circle.  The projection and exterior-connectivity
argument must establish the asserted size bound for the filled interiors.
When the separators are initially winding certificates, fix their
finite family before realizing them in the graph.  Lift a fine
deterministic mesh approximation to the reference path by the
approximate right inverse in the projection argument.  Its projected
path follows the reference path uniformly.  A first exit from a
filled primal circuit forces the reference path to traverse the
corresponding buffered band, up to a vanishing spatial error.
Strict band margins preserve its positive $D$ charge.  Perform the
finite separator-chain construction on the map, take the realization
limit, and only then refine the annulus family.  This avoids assuming
that the projection of a limiting primal circuit is a Jordan curve.

There is also a direct continuum version, used for the tree laws.
Take Euclidean disks as the Jordan domains and assume the stated
boundary connections are continuum certificate connections with
arbitrary open loss.  The same finite chain concatenates at its circle
intersections by the already available continuum joining property.
It has the same cost bound and vanishing endpoint errors.  No primal
circuit reconstruction is needed in this version.  When an upper
comparison has been obtained, local continuity removes its continuum
endpoint errors; alternatively retain open endpoint neighborhoods until
the final extraction, as in the port formulation below.

\begin{corollary}[Comparison from annuli]\label{comp:annulus-comparison}
Assume the hypotheses of Lemma~\ref{comp:circle-chain} hold simultaneously
at every point of every compact subset, along arbitrarily fine families.
If the charged quantity is $D$-length and the boundary connections are
passage connections, then the port upper comparison is $A/a$.
If the charged quantity is certificate cost and the boundary connections
are $D$-connections, then every certificate compactly in $U$ costs at
least $(a/A)D_U(x,y)$.
\end{corollary}

\begin{proof}
The first assertion follows by applying Lemma~\ref{comp:circle-chain} to
a near-minimizing $D_U$-path and then sending $n$ to infinity.  For the
second, perform the construction on discrete realizing paths at each
fixed finite disk family.  Move stopping circles within their open
buffers to preserve strict crossings.  If the number of completed
traversals were unbounded, extract any prescribed finite number of
marked traversals.  The additive cost inequality would give a cost
greater than the total budget once their number exceeds that budget
divided by $a\min_B w_B$.  Thus finitely many marked pieces suffice.
The $D$-connection supplied by the lemma has cost at most $A/a$ times the
budget, and its endpoint errors tend to zero in $D_U$.  This proves the
lower assertion.  Only $D$-continuity removes endpoint errors in this
direction.
\end{proof}

\begin{lemma}[Positive lower comparisons are intrinsic]
\label{comp:intrinsic-lower}
Suppose normalized edge lengths tend to zero, marked cost subdivisions
are available, and $c>0$ is a lower comparison.  A certificate of cost $a$
compactly confined to $U$ satisfies $a\ge cD_U(x,y)$.
\end{lemma}

\begin{proof}
Mark approximating paths at rational fractions of their total costs.
An individual overshoot is at most one edge length and tends to zero.
After simultaneous extraction, marks $x_t$, $t\in\mathbb Q\cap[0,1]$,
satisfy $D(x_s,x_t)\le a|s-t|/c$.  Their Lipschitz extension is a path
of $D$-length at most $a/c$ in the retained compact subset of $U$, with
the required endpoints.  The definition of $D_U$ gives the result.
\end{proof}
